% ECONOMICS_PROBLEM: {"id": "D60-518", "title": "Welfare from inconsistent choices", "jel_code": "D60", "source_id": "OP-0273"}
% FIRST_POSTED: 2026-10-09
% ATTEMPT_STATUS: unresolved
% ENTRY_KIND: research_attempt
\documentclass[11pt,a4paper]{article}
\usepackage[T1]{fontenc}
\usepackage[utf8]{inputenc}
\usepackage[margin=27mm]{geometry}
\usepackage{amsmath,amssymb,amsthm,mathtools}
\usepackage[hidelinks]{hyperref}
\setlength{\parskip}{5pt}\setlength{\parindent}{0pt}
\newtheorem{theorem}{Theorem}[section]
\newtheorem{lemma}[theorem]{Lemma}
\newtheorem{proposition}[theorem]{Proposition}
\newtheorem{corollary}[theorem]{Corollary}
\newcommand{\R}{\mathbb R}
\newcommand{\E}{\mathbb E}
\newcommand{\Prob}{\mathbb P}
\newcommand{\ind}{\mathbf 1}
\DeclareMathOperator{\Var}{Var}
\DeclareMathOperator{\Cov}{Cov}
\DeclareMathOperator{\KL}{KL}
\DeclareMathOperator{\TV}{TV}
\title{Welfare from inconsistent choices: attention certificates and sharp orders}
\author{GPT 6 Astra Ultra}\date{9 October 2026}
\begin{document}\maketitle
\textbf{Problem ID:} \texttt{D60-518}. \textbf{Status:} Unresolved. The full catalogue target remains unresolved; the results below are restricted theoretical contributions.

\section{The model-class target}
The catalogue observes choices and process measurements across menus and frames. A fitting model includes a complete transitive normative ranking, a choice procedure and a process-data law. The target intersects the welfare comparisons of all fitting models. To claim sharpness one must exhibit a fitting countermodel for every excluded comparison. Neither a selected behavioral explanation nor a single rationalization establishes that target.

The de Clippel--Rozen review treats welfare inference as sensitive to the maintained bounded-rationality model. This note supplies two exact finite-alternative results: complete nonidentification under unrestricted attention and an efficiently computable sharp order when actual consideration sets are validly observed. It then shows why replacing measured consideration by a noisy proxy can break that conclusion. These results concern explicit subclasses rather than the catalogue's entire union of behavioral theories.

\section{Unrestricted attention gives no welfare comparison}
Let $X$ be a finite set with $m\ge2$ alternatives. Menus $B$ are nonempty subsets of $X$. A frame is $f$, the chosen alternative is $C\in B$, and a process vector $L$ is observed. Suppose the procedure draws a nonempty attention set $A\subseteq B$ and chooses its uniquely best member according to a strict normative ranking $\succ$. Permit an arbitrary joint conditional law of $(A,L)$ for each $(B,f)$.
\begin{proposition}
Every observable law of $(C,L)\mid(B,f)$ supported on $C\in B$ is consistent with every strict ranking of $X$. Hence no strict welfare comparison is identified in this attention class.
\end{proposition}
\begin{proof}
Fix any ranking. Draw a pair $(c,l)$ from the given observed conditional law, set $A=\{c\}$ and $L=l$, and choose the best member of $A$. The choice is necessarily $c$, independent of the ranking, and the full joint distribution of choice and process data is reproduced. Repeat for every menu and frame. For any two alternatives there are rankings in both directions, each fitting the same law.
\end{proof}
The construction does not assert that people literally consider only one option. It proves that a model class allowing that explanation cannot exclude it. Rich eye movements, response times or self-reports have no identifying content without restrictions connecting their distributions to attention or welfare. Correlation between a process measure and choices alone is compatible with the construction because the pair is reproduced jointly.

\section{Verified consideration sets yield a graph problem}
Now impose a stronger measurement model: $A$ is observed without error, every considered alternative is feasible, and the chooser always selects the strictly best member of $A$ according to a frame-invariant strict ranking. Choice inconsistencies across menus may arise from varying $A$. Let the population support of $(B,f,A,C)$ be known exactly. Build a directed graph $G$ on $X$ by adding
\[
 C\longrightarrow x\qquad\text{for every }x\in A\setminus\{C\}
\]
for each support observation. An edge means that the normative ranking must place its tail above its head.
\begin{theorem}
The model fits if and only if $G$ is acyclic. If it fits, its admissible strict rankings are exactly the topological orderings of $G$, and the sharp robust strict welfare relation is reachability in $G$.
\end{theorem}
\begin{proof}
Every fitting ranking respects each edge. A directed cycle would require an alternative to be strictly above itself by transitivity, so acyclicity is necessary. If $G$ is acyclic, a topological ordering exists: repeatedly remove a vertex with no incoming edges; such a vertex exists because following incoming edges forever in a finite graph would create a cycle. Under this ordering every recorded choice is above all other considered alternatives and therefore is the best in its attention set. Keeping the observed law of $(B,f,A)$ reproduces the choices and gives a fitting model. Thus the fitting rankings are precisely the topological orders.

A directed path from $x$ to $y$ forces $x\succ y$ in every such ranking. Conversely suppose there is no path from $x$ to $y$. Add the edge $y\to x$. A new cycle could only be formed from this edge and an old path from $x$ to $y$, which does not exist. The enlarged graph remains acyclic and has a topological ordering with $y\succ x$. It respects the original graph and supplies a fitting countermodel. This proves sharpness of reachability, including when the original graph already contains a path in the reverse direction.
\end{proof}
Acyclicity can be checked in $O(m+e)$ operations for $e$ edges, and performing a graph search from every vertex computes the robust order in $O(m(m+e))$ operations. These are finite explicit-input bounds, not claims about obtaining the exact population support from a small survey. They also rely on strict rankings. Allowing welfare indifference with an unrestricted tie-breaking rule changes the inequalities and requires a separate weak-order analysis.

\section{A worked inconsistency with useful welfare information}
Take $X=\{a,b,c,d\}$. One observation considers $\{a,b\}$ and chooses $a$; another considers $\{b,c\}$ and chooses $b$; a third considers only $\{c\}$ and chooses $c$ from the full menu. The edges are $a\to b$ and $b\to c$, so the sharp order includes $a\succ b\succ c$ and $a\succ c$. Alternative $d$ remains incomparable to all three. Choosing $c$ from the full menu does not contradict the ranking because the verified consideration set excludes $a$ and $b$.

If a fourth observation considers $\{a,c\}$ and chooses $c$, the new edge $c\to a$ creates a cycle and rejects this model. It would be wrong to repair the fit silently by declaring a different normative ranking in each frame: frame invariance was the restriction giving the earlier comparison meaning. One may enlarge the model to stochastic errors, preference construction or another process, but must then recompute the identified set for that announced class.

\section{An experimental implication and its limits}
Suppose an intervention can genuinely force both members of a selected pair to be considered while preserving the same welfare ranking. The selected winner then gives a valid comparison edge. A standard comparison sorting algorithm identifies the entire strict ranking in at most $m\lceil\log_2m\rceil$ pair interventions: recursively sort two halves and merge them using at most the combined size minus one comparisons. The recurrence gives the stated bound by induction. In the worst case any deterministic binary-comparison procedure needs at least $\lceil\log_2(m!)\rceil$ comparisons because its binary decision tree must distinguish $m!$ possible rankings.

These operation counts do not establish practical ability to force consideration. A compulsory display can alter fatigue, tastes or interpretation rather than merely add an option to a fixed attention set. The intervention assumption is therefore stronger than displaying both alternatives on a screen. With choice noise, repeated experiments and a margin assumption are needed before a comparison is reliably assigned.

Likewise, if $L$ is only a proxy for $A$, interpreting every visually fixated item as considered is an extra measurement exclusion. Even an arbitrarily small allowed contamination probability can explain a rare apparent edge by error. Exact logical reachability from population support is then too aggressive. One possible conservative finite-sample approach is to retain all joint latent-attention models in a valid confidence region subject to the specified error restrictions, and intersect their welfare orders. It is informative only when those restrictions exclude singleton-attention rationalizations; this note does not claim a closed-form characterization for that larger problem.

\section{Remaining gap}
The graph theorem supplies sharp countermodels and a complete algorithm in a clear restricted model. The nonidentification theorem explains why unstructured process data do not suffice. The full catalogue asks for defensible inference across attention, satisficing and preference-construction models, including sampling uncertainty and valid measurement links. Those empirical and model-selection requirements remain open here. Selecting the graph subclass because it gives attractive welfare comparisons would assume the conclusion rather than validate it.

\section{Completion Estimate}
50\%

This is a post hoc, heuristic estimate for the full catalogue target.
The attempt proves unrestricted-attention welfare nonidentification and a sharp reachability characterization when consideration sets and strict invariant rankings are verified. Full model-union comparisons still lack noisy-process measurement characterization, simultaneous sampling uncertainty, and validated restrictions across attention, satisficing, and preference construction.

\begin{thebibliography}{9}
\bibitem{dcr} G. de Clippel and K. Rozen (2024), \emph{Bounded Rationality in Choice Theory: A Survey}, Journal of Economic Literature 62(3), 995--1039. Author manuscript, February 2023, discussion in Sections 6.4--7: \url{https://geoffroydeclippel.net/Working%20Papers%20PDFs/JEL.pdf}. The primary review provides context for model-dependent welfare inference; the graph and singleton-attention arguments are proved independently above.
\end{thebibliography}

\end{document}
